\section{Introduction}

In this chapter, the authors must clearly, orderly, directly, and precisely state the problem that will be investigated. The importance of the Final Degree Project (FDP), the problem to be solved, the objectives to be achieved, any hypotheses (if applicable), and the delimitations of the study should all be expressed. The introduction must include the following subsections:

\subsection{Problem Definition}

This subsection presents the topic and briefly defines the problem, leaving the detailed development of the entire issue for later. It is customary to formulate questions or research inquiries that indicate the orientation of the study.

\subsection{Objectives}

The objectives are closely related to the research problem and describe what is expected to be achieved upon completing the investigation. Objectives are divided into general and specific objectives. To write objectives, it is common to use verbs such as: describe, determine, analyze, explain, assess, etc.

\subsubsection{General Objective}

This defines the ultimate goal of the work in terms of knowledge. In the context of a research project, the general objective describes the desired outcome. It must be specific and concise, expressed in a single sentence. Furthermore, it should begin with a verb in the infinitive form (ending in “-ar”, “-er”, or “-ir” in Spanish, e.g., "to analyze", "to evaluate", "to determine"), and it should be worded in such a way that it encompasses all the specific objectives.

\subsubsection{Specific Objectives}

These are operational in nature. They represent the logical breakdown and sequence of the general objective.

\subsection{Hypotheses}

Depending on the study, there may be one or more hypotheses, or none at all. If it is necessary to include hypotheses in the work, they must correlate with the objectives.

Hypotheses are conjectures, propositions, or speculations offered as tentative answers to the research problem. They are written as affirmations and should be statistically verifiable.

Hypotheses must align with the problem statement, the objectives, and the data analysis. They serve as guides for the research and help organize the ideas.

\subsection{Justification}

The justification should answer several questions, among them: \textit{Why is the research problem important? Why should it be investigated? Who will benefit from the results?} Remember that the justification can be supported by theoretical or knowledge-based, methodological, social, or practical reasons. It is also advisable to back the statement with data, figures, and testimonials.

Grinnell, Williams, and Unrau\cite{grinnell2009} emphasize that the justification should align with the interests of the proposal reviewers and must capture their attention. If you know who they are, you can analyze their motivations and reflect them in the protocol.

The justification presents reasons for relevance in academic or disciplinary, social, and personal dimensions.

In the \textit{academic dimension}, the author must describe the contributions to the field of knowledge expected from the research. These reasons can range from descriptive to analytical, always aiming to convince of the importance of the research and the significance of its findings for the discipline.

In the \textit{social dimension}, the project manager outlines the benefits that society will gain upon completing the research. The project's relevance here is defined by its social impact and the effect of its results on the local environment. Therefore, it is necessary to explicitly frame the object of study within a context close to the researcher.

In the \textit{personal aspect}, the justification outlines the individual motivations that drive the research project. Hence, the researcher must be genuinely committed to and interested in the topic.

The justification of an academic or scientific research project demands the reasoning behind the work itself—something particularly important that legitimizes the research, aimed at producing knowledge in a specific discipline and impacting society as well.

\subsection{Feasibility}

Feasibility regarding the researcher's knowledge and skills, time, location, and budget must be explicitly stated. It must include technical, economic, and management feasibility. Otherwise, the proposal may be considered “academic, professional, or occupational suicide.”

\subsection{Scope Delimitation of the Problem}

In this section, the coverage of the project is descriptively defined. Problem delimitation refers to clearly and objectively defining the object of study and contextualizing it in terms of time and space in which the project will be developed.

The scope of the problem refers to setting the boundaries of the research—that is, concretely describing what was effectively studied in the work and what was not studied.
